% ============================================================
% Section 117: 伦敦方程、迈斯纳效应与超导能隙测量
% ============================================================
\section{固体物理：伦敦方程、迈斯纳效应与超导能隙测量}
\label{sec:london-meissner-gap}

\begin{modelbox}[题目：迈斯纳效应与能隙的实验探测]
(1) 简单超导体的伦敦方程将超导电流和磁矢势联系起来：
\begin{equation*}
j_s=-\frac{n_se^2}{m_ec}A.
\end{equation*}
利用合适的麦克斯韦方程证明如何得出迈斯纳效应。伦敦穿透深度为何？

(2) 多数超导体存在以费米能级为中心的单电子能隙 $\Delta$。下列实验均可确定 $\Delta$ 的存在和大小，请简单讨论其中任意两个：
(a) $T\ll T_c$ 时的比热测量；
(b) 约瑟夫森结的单电子隧道效应；
(c) 薄膜中的微波（$\hbar\omega\ll\Delta$）吸收；
(d) 块状样品中的红外（$\hbar\omega\ll\Delta$）吸收；
(e) $T\ll T_c$ 时的声衰减。
\end{modelbox}

\begin{tipbox}[考点快速分析]
\begin{itemize}
    \item \textbf{伦敦方程}：$\bm{j}_s=-(n_se^2/m_ec)\bm{A}$，取旋度后联立安培定律
    \item \textbf{迈斯纳效应}：$\nabla^2\bm{B}=\bm{B}/\lambda_L^2$，磁场指数衰减，体内 $B\to0$
    \item \textbf{穿透深度}：$\lambda_L=\sqrt{m_ec^2/(4\pi n_se^2)}$（CGS）或 $\sqrt{m_e/(\mu_0n_se^2)}$（SI）
    \item \textbf{能隙测量}：比热（热激发跨越能隙）、隧道效应（输运探测态密度）
\end{itemize}
\end{tipbox}

\begin{derivationbox}[标准解题过程]
\textbf{(1) 迈斯纳效应的证明及伦敦穿透深度}

伦敦方程（高斯单位制）可写为
\begin{equation}
j_s=-\frac{1}{c\Lambda^2}A,\qquad \Lambda^2=\frac{m_e}{n_se^2}.
\end{equation}

取旋度：
\begin{equation}
\nabla\times j_s=-\frac{1}{c\Lambda^2}\nabla\times A=-\frac{1}{c\Lambda^2}B.
\end{equation}

安培环路定律（静态，忽略位移电流）：
\begin{equation}
\nabla\times B=\frac{4\pi}{c}j_s.
\end{equation}

对安培定律取旋度：
\begin{equation}
\nabla\times(\nabla\times B)=\frac{4\pi}{c}\nabla\times j_s.
\end{equation}

利用 $\nabla\times(\nabla\times B)=\nabla(\nabla\cdot B)-\nabla^2B$ 且 $\nabla\cdot B=0$，得
\begin{equation}
-\nabla^2B=\frac{4\pi}{c}\left(-\frac{1}{c\Lambda^2}B\right)=-\frac{4\pi}{c^2\Lambda^2}B,
\end{equation}
即
\begin{equation}
\nabla^2B=\frac{1}{\lambda_L^2}B,\qquad \lambda_L=\sqrt{\frac{c^2\Lambda^2}{4\pi}}=\sqrt{\frac{m_ec^2}{4\pi n_se^2}}.
\end{equation}

在 SI 单位制中
\begin{equation}
\boxed{\lambda_L=\sqrt{\frac{m_e}{\mu_0n_se^2}}.}
\end{equation}

\textbf{迈斯纳效应}：考虑超导体占据半无限空间 $x>0$，磁场 $B=B(x)\hat{z}$ 平行于表面。方程化为
\begin{equation}
\frac{d^2B}{dx^2}=\frac{1}{\lambda_L^2}B,
\end{equation}
其物理解为指数衰减
\begin{equation}
\boxed{B(x)=B_0e^{-x/\lambda_L}.}
\end{equation}

磁场从表面向体内指数衰减，$x\gg\lambda_L$ 时 $B\to0$。这表明超导体内部磁场为零——\textbf{完全抗磁性（迈斯纳效应）}。$\lambda_L$ 为磁场穿透的特征深度，典型值 $\lambda_L\sim10^{-8}\sim10^{-7}\,\text{m}$（$100\sim1000\,\text{\AA}$）。

\textbf{(2) 超导能隙 $\Delta$ 的实验测定}

\textbf{方法一：$T\ll T_c$ 时的比热测量}

低温下，超导体中的准粒子（拆散的库珀对）需跨越能隙 $2\Delta$ 才能被热激发。根据 BCS 理论，超导态电子比热呈指数衰减：
\begin{equation}
c_s\propto\exp\!\left(-\frac{\Delta}{k_BT}\right).
\end{equation}

实验方法：测量 $T\ll T_c$（如 $T<T_c/3$）下的比热 $c_s(T)$，作 $\ln c_s$ 对 $1/T$ 图，斜率即为 $-\Delta/k_B$，由此精确获得能隙 $\Delta$。典型结果：$2\Delta(0)\approx3.5k_BT_c$。

\textbf{方法二：单电子隧道效应（Giaever 隧道）}

在超导体--绝缘体--正常金属（S--I--N）隧道结中，正常金属电子可隧穿进入超导体。由于超导体在费米能级附近存在能隙 $\pm\Delta$，当偏压 $V<\Delta/e$ 时，正常金属中的电子无法找到可占据的空态（能隙内无能级），隧道电流极小。当 $V=\Delta/e$ 时，电流突然上升。

实验上测量 S--I--N 结的 $I$--$V$ 特性曲线，微分电导 $dI/dV$ 在 $V=\pm\Delta/e$ 处出现尖锐峰，直接给出能隙 $\Delta$。这是测量超导能隙最直接、最精确的方法。
\end{derivationbox}

\begin{formulabox}[答案汇总]
\begin{center}
\begin{tabular}{ll}
\toprule
\textbf{项目} & \textbf{结果} \\
\midrule
(1) 迈斯纳效应 & 伦敦方程 + 麦克斯韦方程导出 $\nabla^2B=B/\lambda_L^2$，指数衰减 \\
(1) 伦敦穿透深度 & $\displaystyle\lambda_L=\sqrt{\frac{m_ec^2}{4\pi n_se^2}}$（CGS）或 $\sqrt{m_e/(\mu_0n_se^2)}$（SI） \\
(2) 比热测量 & 低温 $c_s\propto\exp(-\Delta/k_BT)$，指数拟合斜率得 $\Delta$ \\
(2) 单电子隧道 & S--I--N 结 $I$--$V$ 曲线在 $V=\Delta/e$ 处电流陡升 \\
\bottomrule
\end{tabular}
\end{center}
\end{formulabox}

\begin{insightbox}[物理图像]
\begin{itemize}
    \item 伦敦方程是描述超导体电磁响应的唯象方程，成功预言了迈斯纳效应和穿透深度。迈斯纳效应是超导体的根本特征之一，与零电阻共同构成超导态的完整定义。
    \item 多种独立的实验方法（比热、隧道、红外吸收、超声衰减）对能隙的测定结果一致，验证了 BCS 理论的正确性。比热和隧道效应是最经典的两种方法，分别从热力学和输运角度提供能隙证据。
\end{itemize}
\end{insightbox}
